\documentclass[10pt]{article}
\usepackage[margin=1in]{geometry}
\usepackage[T1]{fontenc}
\usepackage{lmodern,microtype,amsmath,amssymb,amsthm}
\theoremstyle{definition}
\newtheorem{assumption}{Assumption}
\begin{document}
\section*{Stationarity Under Strong Growth}
We consider minimizing a possibly nonconvex, twice continuously differentiable function $f:E\to\mathbb{R}$ on a finite-dimensional real inner product space $E$. Its gradient and Hessian are Lipschitz with constants $L>0$ and $\rho>0$, and a given point $x^0\in E$ satisfies $f(x^0)-\inf_{x\in E}f(x)\leq\Delta$, where $\Delta\geq0$. The oracle returns unbiased stochastic gradients with independent, identically distributed samples and mean squared error at most $(\kappa-1)\|\nabla f(x)\|^2$, where $\kappa\geq1$. The goal is a measurable output $\bar{x}\in E$ with $\mathbb{E}\|\nabla f(\bar{x})\|\leq\varepsilon$, where $\varepsilon>0$ satisfies $4^8\rho\varepsilon\leq L^2$.

The optimization problem is
\[
\inf_{x\in E} f(x).
\]
The norm $\|\cdot\|$ is induced by the inner product on $E$, and $\|\cdot\|_{\mathrm{op}}$ denotes the induced operator norm. The algorithm receives $E$, the measurable sample space $\Omega$, $L$, $\rho$, $\kappa$, $\Delta$, $\varepsilon$, and $x^0$.

Our assumptions:
\begin{assumption}[Gradient and Hessian regularity]
The function $f$ belongs to $C^2(E;\mathbb{R})$, and, for all $x,y\in E$,
\[
\begin{aligned}
\|\nabla f(x)-\nabla f(y)\|&\leq L\|x-y\|,\\
\|\nabla^2 f(x)-\nabla^2 f(y)\|_{\mathrm{op}}&\leq\rho\|x-y\|.
\end{aligned}
\]
Here $\nabla^2 f(x):E\to E$ is the derivative of the gradient.
\end{assumption}

\begin{assumption}[Initial objective gap]
The given point $x^0$ satisfies
\[
f(x^0)-f(x)\leq\Delta,\qquad x\in E.
\]
\end{assumption}

\begin{assumption}[Strong-growth stochastic oracle]
There is a probability measure $\nu$ on $\Omega$ and a jointly measurable map $g:E\times\Omega\to E$. For every $x\in E$, the random vector $g(x,\xi)$ is Bochner integrable and $\|g(x,\xi)-\nabla f(x)\|^2$ is integrable for $\xi\sim\nu$, with
\[
\begin{aligned}
\mathbb{E}_{\xi\sim\nu}[g(x,\xi)]&=\nabla f(x),\\
\mathbb{E}_{\xi\sim\nu}\bigl[\|g(x,\xi)-\nabla f(x)\|^2\bigr]
&\leq(\kappa-1)\|\nabla f(x)\|^2.
\end{aligned}
\]
The samples $(\xi_k)_{k\geq0}$ are measurable, mutually independent, and each has law $\nu$ under a joint probability law $Q$ on $\Omega^{\mathbb{N}}$. At the $k$th query $x\in E$, the oracle returns $g(x,\xi_k)$. Algorithm randomness is independent of this sample sequence.
\end{assumption}

Queries may be made at any point of $E$. After the allotted oracle calls, the output $\bar{x}$ must be measurable, its gradient norm must be integrable, and it must satisfy
\[
\mathbb{E}_{P}\bigl[\|\nabla f(\bar{x})\|\bigr]\leq\varepsilon,
\]
where $P$ is the joint law of the algorithm's randomness and the oracle samples.
\end{document}
